\documentclass[12pt,a4paper]{article}
\usepackage[T1]{fontenc}
\usepackage{amsmath,amssymb,mathrsfs,tikz,times,pifont}
\usepackage{enumitem}
\usepackage{multicol}
\usepackage{lmodern}
\usetikzlibrary{trees}
\newcommand\circitem[1]{%
\tikz[baseline=(char.base)]{
\node[circle,draw=gray, fill=red!55,
minimum size=1.2em,inner sep=0] (char) {#1};}}
\newcommand\boxitem[1]{%
\tikz[baseline=(char.base)]{
\node[fill=cyan,
minimum size=1.2em,inner sep=0] (char) {#1};}}
\setlist[enumerate,1]{label=\protect\circitem{\arabic*}}
\setlist[enumerate,2]{label=\protect\boxitem{\alph*}}
%%%::::::by chnini ameur :::::::%%%
\everymath{\displaystyle}
\usepackage[left=1cm,right=1cm,top=1cm,bottom=1.7cm]{geometry}
\usepackage[colorlinks=true, linkcolor=red, urlcolor=red, citecolor=red]{hyperref}
\usepackage{array,multirow}
\usepackage[most]{tcolorbox}
\usepackage{varwidth}
\usepackage{float} %pour utiliser l'option [H] qui force l'image à apparaître exactement à l'endroit où elle est placée dans le code.
\tcbuselibrary{skins,hooks}
\usetikzlibrary{patterns}
%%%::::::by chnini ameur :::::::%%%
\newtcolorbox{exa}[2][]{enhanced,breakable,before skip=2mm,after skip=5mm,
colback=yellow!20!white,colframe=black!20!blue,boxrule=0.5mm,
attach boxed title to top left ={xshift=0.6cm,yshift*=1mm-\tcboxedtitleheight},
fonttitle=\bfseries,
title={#2},#1,
% varwidth boxed title*=-3cm,
boxed title style={frame code={
\path[fill=tcbcolback!30!black]
([yshift=-1mm,xshift=-1mm]frame.north west)
arc[start angle=0,end angle=180,radius=1mm]
([yshift=-1mm,xshift=1mm]frame.north east)
arc[start angle=180,end angle=0,radius=1mm];
\path[left color=tcbcolback!60!black,right color = tcbcolback!60!black,
middle color = tcbcolback!80!black]
([xshift=-2mm]frame.north west) -- ([xshift=2mm]frame.north east)
[rounded corners=1mm]-- ([xshift=1mm,yshift=-1mm]frame.north east)
-- (frame.south east) -- (frame.south west)
-- ([xshift=-1mm,yshift=-1mm]frame.north west)
[sharp corners]-- cycle;
},interior engine=empty,
},interior style={top color=yellow!5}}
%%%%%%%%%%%%%%%%%%%%%%%
\usepackage{fancyhdr}
\usepackage{eso-pic}         % Pour ajouter des éléments en arrière-plan
% Commande pour ajouter du texte en arrière-plan
\usepackage{tkz-tab}
% Filigrane "PGB" centré sur chaque page
\AddToShipoutPicture{
  \AtTextCenter{%
    \makebox[0pt]{\rotatebox{50}{\textcolor[gray]{0.6}{\fontsize{1.5cm}{1.5cm}\selectfont maths-mes-maths}}}
  }
}
\usepackage{lastpage}
\fancyhf{}
\pagestyle{fancy}
\renewcommand{\footrulewidth}{1pt}
\renewcommand{\headrulewidth}{0pt}
\renewcommand{\footruleskip}{10pt}
\fancyfoot[R]{
\color{red}\ding{45}\ \textbf{2026}
}
\fancyfoot[L]{
\color{red}\ding{45}\ \textbf{Prof:M. Diatta}
}
\cfoot{\bf
\thepage /
\pageref{LastPage}}
% Création du compteur pour les exercices
\newcounter{exercice}
\renewcommand{\theexercice}{\arabic{exercice}}  % Définit l'affichage du compteur en chiffres arabes

% Définir la commande \exo
\newcommand{\exo}{\refstepcounter{exercice}\textbf{Exercice \theexercice} }
\begin{document}
\renewcommand{\arraystretch}{1.5}
\renewcommand{\arrayrulewidth}{1.2pt}
\begin{tcolorbox}[
    colback=white,
    colframe=blue!60!black,
    boxrule=0.8pt,
    arc=3pt,
    left=3mm, right=3mm, top=2.5mm, bottom=2.5mm
]
    \begin{center}
        \textbf{\large Lycée de Diattacounda} \hfill \textbf{A.S. : 2026/2027} \\[1mm]
        \textbf{Mathématiques} $\bullet$ \textbf{Niveau : TS2} \hfill \textbf{Date : 13/11/2026} \\[2mm]
        \rule{0.8\linewidth}{0.4pt} \\[1.5mm]
        {\Large\bfseries\sffamily TD : Nombres Complexes}
    \end{center}
\end{tcolorbox}
\vspace{4mm}


\fbox{\textbf{\exo}} [Forme algébrique et calcul de $i^n$]

\begin{enumerate}
    \item Formes algébriques des nombres complexes

\textbf{Rappel :} La forme algébrique d'un nombre complexe s'écrit $z = a + ib$, avec $a, b \in \mathbb{R}$ et $i^2 = -1$.

\begin{enumerate}
    \item $z_1 = i(1 - 4i) + 3(2 - i)$
    
$    
\begin{aligned}
        z_1 &= i - 4i^2 + 6 - 3i \\
        &= i + 4 + 6 - 3i \\
        &= 10 - 2i
\end{aligned}
$

    C'est un nombre complexe quelconque (ni réel, ni imaginaire pur).

    \item $z_2 = (3 - i)^2 + 6i$
    
    $    
\begin{aligned}
        z_2 &= 3^2 - 2(3)(i) + i^2 + 6i \\
        &= 9 - 6i - 1 + 6i \\
        &= 8
\end{aligned}
$

    C'est un nombre \textbf{réel}.

    \item $z_3 = \dfrac{3}{1 - i}$
    
    $    
\begin{aligned}
        z_3 &= \frac{3(1 + i)}{(1 - i)(1 + i)} \\
        &= \frac{3 + 3i}{1^2 - i^2} \\
        &= \frac{3 + 3i}{2} \\
        &= \frac{3}{2} + \frac{3}{2}i
\end{aligned}
$

    C'est un nombre complexe quelconque.

    \item $z_4 = \dfrac{1 + 2i}{2 - i}$

    $    
\begin{aligned}
        z_4 &= \frac{(1 + 2i)(2 + i)}{(2 - i)(2 + i)} \\
        &= \frac{2 + i + 4i + 2i^2}{2^2 - i^2} \\
        &= \frac{2 + 5i - 2}{4 + 1} \\
        &= \frac{5i}{5} \\
        &= i
\end{aligned}
$

    C'est un nombre \textbf{imaginaire pur}.

    \item $z_5 = (2 + i\sqrt{3})(\sqrt{3} - i)$

    $    
\begin{aligned}
        z_5 &= 2\sqrt{3} - 2i + i\sqrt{3}\cdot\sqrt{3} - i^2\sqrt{3} \\
        &= 2\sqrt{3} - 2i + 3i + \sqrt{3} \\
        &= 3\sqrt{3} + i
\end{aligned}
$

    C'est un nombre complexe quelconque.

    \item $z_6 = (2 + 5i)^2$

    $    
\begin{aligned}
        z_6 &= 2^2 + 2(2)(5i) + (5i)^2 \\
        &= 4 + 20i + 25i^2 \\
        &= 4 + 20i - 25 \\
        &= -21 + 20i
\end{aligned}
$

    C'est un nombre complexe quelconque.

    \item $z_7 = \dfrac{3 - 3i}{1 + i}$

    $    
\begin{aligned}
        z_7 &= \frac{(3 - 3i)(1 - i)}{(1 + i)(1 - i)} \\
        &= \frac{3 - 3i - 3i + 3i^2}{1^2 - i^2} \\
        &= \frac{3 - 6i - 3}{2} \\
        &= \frac{-6i}{2} \\
        &= -3i
\end{aligned}
$

    C'est un nombre \textbf{imaginaire pur}.
\end{enumerate}

\item Calcul des puissances de $i$

\textbf{Rappel :} $i^0 = 1$, $i^1 = i$, $i^2 = -1$, $i^3 = -i$ et $i^4 = 1$. Si $n = 4q + r$ (avec $0 \le r < 4$), alors $i^n = i^r$.

\begin{enumerate}
    \item $2025 = 4 \times 506 + 1 \implies i^{2025} = i^{4 \times 506 + 1} =(i^4)^{506} \times i^1 = 1^{506} \times i = \mathbf{i}$
    \item $40 = 4 \times 10 + 0 \implies i^{40} = i^{4 \times 10} =(i^4)^{10} = 1^{10} = \mathbf{1}$
    \item $27 = 4 \times 6 + 3 \implies i^{27} = i^{4 \times 6 + 3} = (i^4)^6 \times i^3 = 1^6 \times (-i) = \mathbf{-i}$
    \item $34 = 4 \times 8 + 2 \implies i^{34} = i^{4 \times 8 + 2} =(i^4)^8 \times i^2 = 1^8 \times (-1) = \mathbf{-1}$
\end{enumerate}

\end{enumerate}

\fbox{\textbf{\exo}} [Conjugué]

\textbf{Rappel :} Si $z = a + ib$ (avec $a, b \in \mathbb{R}$), alors son conjugué est $\overline{z} = a - ib$. \\
On utilise les propriétés suivantes : $\overline{z + z'} = \overline{z} + \overline{z'}$, $\overline{z \times z'} = \overline{z} \times \overline{z'}$, $\overline{z^n} = (\overline{z})^n$ et $\overline{\left(\dfrac{z}{z'}\right)} = \dfrac{\overline{z}}{\overline{z'}}$ (avec $z' \neq 0$).

\begin{enumerate}
    \item \textbf{Déterminer le conjugué des nombres complexes suivants :}
    \begin{enumerate}
        \item $z_1 = 3i$
        \[
            \overline{z_1} = \overline{3i} = \mathbf{-3i}
        \]

        \item $z_2 = -i(1 + 4i)$
        
        $\begin{aligned}
            \overline{z_2} &= \overline{-i} \times \overline{(1 + 4i)} \\
            &= i(1 - 4i) \\
            &= i - 4i^2 \\
            &= \mathbf{4 + i}
        \end{aligned}$

        \item $z_3 = (3 - 5i)^3$
        
        $\begin{aligned}
            \overline{z_3} &= \overline{(3 - 5i)^3} \\
            &= \left(\overline{3 - 5i}\right)^3 \\
            &= \mathbf{(3 + 5i)^3}
        \end{aligned}$
        
        \textit{(Sous forme algébrique développée : $\overline{z_3} = 27 + 135i - 225 - 125i = \mathbf{-198 + 10i}$.)}

        \item $z_4 = \dfrac{2i}{i + 3} - \dfrac{3}{2 + i}$
        
        $\begin{aligned}
            \overline{z_4} &= \overline{\left(\frac{2i}{3 + i}\right)} - \overline{\left(\frac{3}{2 + i}\right)} \\
            &= \frac{\overline{2i}}{\overline{3 + i}} - \frac{\overline{3}}{\overline{2 + i}} \\
            &= \mathbf{\frac{-2i}{3 - i} - \frac{3}{2 - i}}
        \end{aligned}$
        
        \textit{Mise sous forme algébrique :}
        
        $\begin{aligned}
            \overline{z_4} &= \frac{-2i(3 + i)}{(3 - i)(3 + i)} - \frac{3(2 + i)}{(2 - i)(2 + i)} \\
            &= \frac{-6i - 2i^2}{9 + 1} - \frac{6 + 3i}{4 + 1} \\
            &= \frac{2 - 6i}{10} - \frac{6 + 3i}{5} \\
            &= \frac{1 - 3i}{5} - \frac{6 + 3i}{5} \\
            &= \mathbf{-1 - \frac{6}{5}i}
        \end{aligned}$
        
    \end{enumerate}

    \item \textbf{Étude de $z_5$ et $z_6$ :}
    \begin{itemize}
        \item \textbf{Démonstration :}
        
        $\begin{aligned}
            \overline{z_6} &= \overline{\left(\frac{2 + i}{1 - i}\right)} \\
            &= \frac{\overline{2 + i}}{\overline{1 - i}} \\
            &= \frac{2 - i}{1 + i} \\
            &= z_5
        \end{aligned}$
        
        On a donc bien $\overline{z_6} = z_5$ (soit $z_6 = \overline{z_5}$).

        \item \textbf{Déduction :} \\
        Pour tout nombre complexe $z$, $z + \overline{z} = 2\operatorname{Re}(z) \in \mathbb{R}$ et $z - \overline{z} = 2i\operatorname{Im}(z) \in i\mathbb{R}$. \\
        Comme $z_6 = \overline{z_5}$, on en déduit sans calcul que :
        \begin{itemize}
            \item $z_5 + z_6 = z_5 + \overline{z_5}$ est un \textbf{nombre réel}.
            \item $z_5 - z_6 = z_5 - \overline{z_5}$ est un \textbf{nombre imaginaire pur}.
        \end{itemize}
    \end{itemize}
\end{enumerate}

\fbox{\textbf{\exo}} [Module et argument]

\textbf{Rappels :} Pour $z = a + ib$ (avec $a, b \in \mathbb{R}$) non nul :
\begin{itemize}
    \item Le module est $|z| = \sqrt{a^2 + b^2}$.
    \item Un argument $\theta = \arg(z)$ vérifie $\cos(\theta) = \dfrac{a}{|z|}$ et $\sin(\theta) = \dfrac{b}{|z|}$.
    \item Propriétés des modules et arguments : $|z \cdot z'| = |z| \cdot |z'|$, $\left|\dfrac{z}{z'}\right| = \dfrac{|z|}{|z'|}$, $|z^n| = |z|^n$, $\arg(z \cdot z') \equiv \arg(z) + \arg(z') \pmod{2\pi}$, $\arg\left(\dfrac{z}{z'}\right) \equiv \arg(z) - \arg(z') \pmod{2\pi}$ et $\arg(z^n) \equiv n\arg(z) \pmod{2\pi}$.
\end{itemize}

\begin{enumerate}
    \item $z_1 = -\sqrt{3} + i$
    
    $
    \begin{aligned}
        |z_1| &= \sqrt{(-\sqrt{3})^2 + 1^2} 
        &= \sqrt{3 + 1}\\ 
        &= \mathbf{2} 
    \end{aligned}
    $
    
    $
    \begin{aligned}
        z_1 &= 2\left(-\frac{\sqrt{3}}{2} + \frac{1}{2}i\right)\\
        &= 2\left(\cos\frac{5\pi}{6} + i\sin\frac{5\pi}{6}\right) \implies \arg(z_1) \equiv \mathbf{\frac{5\pi}{6}} \pmod{2\pi}
    \end{aligned}
    $
    
    \item $z_2 = -1 + i\sqrt{3}$
    
    $
    \begin{aligned}
        |z_2| &= \sqrt{(-1)^2 + (\sqrt{3})^2}\\ 
        &= \sqrt{1 + 3}\\ 
        &= \mathbf{2} 
    \end{aligned}$

	$
    \begin{aligned}
        z_2 &= 2\left(-\frac{1}{2} + \frac{\sqrt{3}}{2}i\right)\\ 
        &= 2\left(\cos\frac{2\pi}{3} + i\sin\frac{2\pi}{3}\right) \implies \arg(z_2) \equiv \mathbf{\frac{2\pi}{3}} \pmod{2\pi}
    \end{aligned}$

    \item $z_3 = 1 - i$
    
    $
    \begin{aligned}
        |z_3| &= \sqrt{1^2 + (-1)^2}\\
        &= \mathbf{\sqrt{2}} \\
    \end{aligned}$

    $\begin{aligned}
        z_3 &= \sqrt{2}\left(\frac{\sqrt{2}}{2} - \frac{\sqrt{2}}{2}i\right)\\
        &= \sqrt{2}\left(\cos\left(-\frac{\pi}{4}\right) + i\sin\left(-\frac{\pi}{4}\right)\right) \implies \arg(z_3) \equiv \mathbf{-\frac{\pi}{4}} \pmod{2\pi}
    \end{aligned}$
    
    \item $z_4 = 1 + i$
    
    $
    \begin{aligned}
        |z_4| &= \sqrt{1^2 + 1^2}\\
        &= \mathbf{\sqrt{2}}
    \end{aligned}
    $
    
    $
    \begin{aligned}
        z_4 &= \sqrt{2}\left(\frac{\sqrt{2}}{2} + \frac{\sqrt{2}}{2}i\right)\\
        &= \sqrt{2}\left(\cos\frac{\pi}{4} + i\sin\frac{\pi}{4}\right) \implies \arg(z_4) \equiv \mathbf{\frac{\pi}{4}} \pmod{2\pi}
    \end{aligned}
    $
    
    \item $z_5 = -\sqrt{3} + i$ \\
    Identique à $z_1$ : $\mathbf{|z_5| = 2}$ et $\mathbf{\arg(z_5) \equiv \dfrac{5\pi}{6} \pmod{2\pi}}$.

	$
    \begin{aligned}
        |z_5| &= \sqrt{(-\sqrt{3})^2 + 1^2}\\
        &= \sqrt{3 + 1}\\
        &= \mathbf{2}
    \end{aligned}
    $
    
    $
    \begin{aligned}
        z_5 &= 2\left(-\frac{\sqrt{3}}{2} + \frac{1}{2}i\right)\\
        &= 2\left(\cos\frac{5\pi}{6} + i\sin\frac{5\pi}{6}\right) \implies \arg(z_5) \equiv \mathbf{\frac{5\pi}{6}} \pmod{2\pi}
    \end{aligned}
    $

    \item .
    
	$
    \begin{aligned}
    z_6 &= \dfrac{2i - 2\sqrt{3}}{4i + 4}\\ 
    &= \dfrac{2(-\sqrt{3} + i)}{4(1 + i)}\\ 
    &= \dfrac{1}{2} \cdot \dfrac{z_1}{z_4}
        \end{aligned}
    $
    
    $
    \begin{aligned}
        |z_6| &= \frac{1}{2} \cdot \frac{|z_1|}{|z_4|}\\
        &= \frac{1}{2} \cdot \frac{2}{\sqrt{2}}\\
        &= \mathbf{\frac{\sqrt{2}}{2}}
    \end{aligned}
    $
    
    $
    \begin{aligned}
        \arg(z_6) &\equiv \arg(z_1) - \arg(z_4) \pmod{2\pi}\\
        &\equiv \frac{5\pi}{6} - \frac{\pi}{4} \pmod{2\pi}\\
        &\equiv \frac{10\pi - 3\pi}{12}\\ 
        &\equiv \mathbf{\frac{7\pi}{12}} \pmod{2\pi}
    \end{aligned}
    $

    \item $z_7 = (\sqrt{3} + i)(1 - i)^2$
    
    On pose $a = \sqrt{3} + i$, d'où $|a| = 2$ et $\arg(a) \equiv \dfrac{\pi}{6} \pmod{2\pi}$. 
    
    On a $z_7 = a \cdot z_3^2$ :
    
    $
    \begin{aligned}
        |z_7| &= |a| \cdot |z_3|^2\\
        &= 2 \cdot (\sqrt{2})^2\\
        &= \mathbf{4}
    \end{aligned}
    $
    
    $
    \begin{aligned}
        \arg(z_7) &\equiv \arg(a) + 2\arg(z_3) \pmod{2\pi}\\
        &\equiv \frac{\pi}{6} + 2\left(-\frac{\pi}{4}\right) \pmod{2\pi}\\
        &\equiv \frac{\pi}{6} - \frac{\pi}{2}\\ 
        &\equiv \mathbf{-\frac{\pi}{3}} \pmod{2\pi}
    \end{aligned}
    $
    
    \item $z_8 = \dfrac{(-\sqrt{3} + i)^3}{(1 + i)^2} = \dfrac{z_1^3}{z_4^2}$
    
    $
    \begin{aligned}
        |z_8| &= \frac{|z_1|^3}{|z_4|^2}\\
        &= \frac{2^3}{(\sqrt{2})^2}\\
        &= \frac{8}{2}\\ 
        &= \mathbf{4}
    \end{aligned}
    $
    
    $
    \begin{aligned}
        \arg(z_8) &\equiv 3\arg(z_1) - 2\arg(z_4) \pmod{2\pi}\\
        &\equiv 3\left(\frac{5\pi}{6}\right) - 2\left(\frac{\pi}{4}\right) \pmod{2\pi}\\
        &\equiv \frac{5\pi}{2} - \frac{\pi}{2}\\ 
        &\equiv 2\pi\\ 
        &\equiv \mathbf{0} \pmod{2\pi}
    \end{aligned}
    $

    \item $z_9 = 2 \dfrac{(\sqrt{3} + i)^2}{(-1 - i)^2}$
    
    On a $|\sqrt{3} + i| = 2$ avec $\arg(\sqrt{3} + i) \equiv \dfrac{\pi}{6} \pmod{2\pi}$, et $|-1 - i| = \sqrt{2}$ 
    
    avec $\arg(-1 - i) \equiv -\dfrac{3\pi}{4} \pmod{2\pi}$.
    
    $
    \begin{aligned}
        |z_9| &= 2 \cdot \frac{2^2}{(\sqrt{2})^2}\\
        &= 2 \cdot \frac{4}{2}\\
        &= \mathbf{4}
    \end{aligned}
    $
    
    $
    \begin{aligned}
        \arg(z_9) &\equiv \arg(2) + 2\arg(\sqrt{3} + i) - 2\arg(-1 - i) \pmod{2\pi}\\
        &\equiv 0 + 2\left(\frac{\pi}{6}\right) - 2\left(-\frac{3\pi}{4}\right) \pmod{2\pi}\\
        &\equiv \frac{\pi}{3} + \frac{3\pi}{2} \\
        &\equiv \frac{11\pi}{6} \\
        &\equiv \mathbf{-\frac{\pi}{6}} \pmod{2\pi}
    \end{aligned}
    $

    \item $z_{10} = -2i (1 + i\sqrt{3})^6$
    
    On a $|-2i| = 2$ avec $\arg(-2i) \equiv -\dfrac{\pi}{2} \pmod{2\pi}$, et $|1 + i\sqrt{3}| = 2$ avec $\arg(1 + i\sqrt{3}) \equiv \dfrac{\pi}{3} \pmod{2\pi}$.
    
    $
    \begin{aligned}
        |z_{10}| &= |-2i| \cdot |1 + i\sqrt{3}|^6\\
        &= 2 \cdot 2^6\\
        &= 2^7 \\
        &= \mathbf{128}
    \end{aligned}
    $
    
    $
    \begin{aligned}
        \arg(z_{10}) &\equiv \arg(-2i) + 6\arg(1 + i\sqrt{3}) \pmod{2\pi}\\
        &\equiv -\frac{\pi}{2} + 6\left(\frac{\pi}{3}\right) \pmod{2\pi}\\
        &\equiv -\frac{\pi}{2} + 2\pi \\
        &\equiv \mathbf{-\frac{\pi}{2}} \pmod{2\pi}
    \end{aligned}
    $
\end{enumerate}

\fbox{\textbf{\exo}} [Formes trigonométrique et exponentielle]

\fbox{\textbf{\exo}} [Forme trigonométrique]

\fbox{\textbf{\exo}} [Classique !]

\fbox{\textbf{\exo}} [Racines carrées et Equations]

\fbox{\textbf{\exo}} [Racines $n$-ièmes]

On donne $z_0 = 1 - i\sqrt{3}$

\begin{enumerate}
    \item Donner une écriture trigonométrique de $z_0$.
    \item Montrer que $z_0^4 = -8 + 8i\sqrt{3}$.
    \item Résoudre dans $\mathbb{C}$ l'équation $z^4 = 1$.
    \item En déduire les solutions de $(E) : z^4 = -8 + 8i\sqrt{3}$ sous forme algébrique et sous forme trigonométrique.
\end{enumerate}
\fbox{\textbf{\exo}} [Racines $n$-ièmes]

\fbox{\textbf{\exo}}[Linéarisation]

\fbox{\textbf{\exo}} [Lieux géométriques]

\fbox{\textbf{\exo}} [Lieux géométriques]

\fbox{\textbf{\exo}} [Deux méthodes pour un problème]

\fbox{\textbf{\exo}} [Complexe et géométrie]

\fbox{\textbf{\exo}}

\fbox{\textbf{\exo}}

\fbox{\textbf{\exo}}

\fbox{\textbf{\exo}}

\fbox{\textbf{\exo}}

\fbox{\textbf{\exo}}

\fbox{\textbf{\exo}}

\fbox{\textbf{\exo}}

\fbox{\textbf{\exo}}

\fbox{\textbf{\exo}}

\fbox{\textbf{\exo}}

\fbox{\textbf{\exo}}
\section*{\underline{\textbf{Exercice :}} Nombres complexes}

On donne 
\[
a = \sqrt{2 - \sqrt{3}} - i\sqrt{2 + \sqrt{3}}.
\]

\begin{enumerate}
    \item \textbf{Calcul de $a^2$ et du module de $a$}

    Posons 
    \[
    x = \sqrt{2 - \sqrt{3}}, \qquad y = \sqrt{2 + \sqrt{3}}.
    \]
    Alors $a = x - iy$.

    On calcule :
    \[
    a^2 = (x - iy)^2 = x^2 - 2ixy - y^2.
    \]

    Or
    \[
    x^2 = 2 - \sqrt{3}, \qquad y^2 = 2 + \sqrt{3}.
    \]

    Donc
    \[
    a^2 = (2 - \sqrt{3}) - (2 + \sqrt{3}) - 2i\sqrt{(2-\sqrt3)(2+\sqrt3)}
    = -2\sqrt{3} - 2i.
    \]

    Ainsi,
    \[
    \boxed{a^2 = -2\sqrt{3} - 2i.}
    \]

    Le module de $a$ vérifie :
    \[
    |a|^2 = |a^2| = \sqrt{(-2\sqrt3)^2 + (-2)^2} = \sqrt{12 + 4} = 4,
    \]
    d'où
    \[
    \boxed{|a| = 2.}
    \]

    \item \textbf{Forme trigonométrique de $a^2$ et argument de $a$}

    On a
    \[
    a^2 = -2\sqrt{3} - 2i.
    \]

    Son module est $|a^2| = 4$. On écrit
    \[
    a^2 = 4(\cos\theta + i\sin\theta),
    \]
    avec
    \[
    \cos\theta = \frac{-2\sqrt3}{4} = -\frac{\sqrt3}{2}, 
    \qquad
    \sin\theta = \frac{-2}{4} = -\frac{1}{2}.
    \]

    On reconnaît
    \[
    \theta = \frac{7\pi}{6} \quad (\text{ou } -\tfrac{5\pi}{6}).
    \]

    Donc
    \[
    \boxed{a^2 = 4\left(\cos\frac{7\pi}{6} + i\sin\frac{7\pi}{6}\right).}
    \]

    Comme $a^2 = |a|^2 e^{i2\arg(a)}$, on a :
    \[
    2\arg(a) = \frac{7\pi}{6} + 2k\pi
    \quad\Rightarrow\quad
    \arg(a) = \frac{7\pi}{12} + k\pi.
    \]

    Une valeur positive possible est donc :
    \[
    \boxed{\arg(a) = \frac{19\pi}{12}.}
    \]

    \item \textbf{Valeurs exactes des cosinus et sinus}

    On a 
    \[
    a = 2\left(\cos\frac{19\pi}{12} + i\sin\frac{19\pi}{12}\right).
    \]

    Donc
    \[
    \cos\frac{19\pi}{12} = \frac{\Re(a)}{2} = \frac{\sqrt{2-\sqrt3}}{2}, 
    \]
    \[
    \sin\frac{19\pi}{12} = \frac{\Im(a)}{2} = -\frac{\sqrt{2+\sqrt3}}{2}.
    \]

    Or 
    \[
    \frac{19\pi}{12} = \frac{7\pi}{12} + \pi,
    \]
    donc
    \[
    \cos\frac{19\pi}{12} = -\cos\frac{7\pi}{12}, 
    \qquad
    \sin\frac{19\pi}{12} = -\sin\frac{7\pi}{12}.
    \]

    Ainsi :
    \[
    \boxed{
    \begin{aligned}
    \cos\frac{7\pi}{12} &= \frac{\sqrt{2-\sqrt3}}{2},\\
    \sin\frac{7\pi}{12} &= \frac{\sqrt{2+\sqrt3}}{2}.
    \end{aligned}}
    \]

    Enfin,
    \[
    \boxed{
    \begin{aligned}
    \cos\frac{\pi}{12} &= \frac{\sqrt{2+\sqrt3}}{2},\\
    \sin\frac{\pi}{12} &= \frac{\sqrt{2-\sqrt3}}{2}.
    \end{aligned}}
    \]

\end{enumerate}

\fbox{\textbf{\exo}}

\fbox{\textbf{\exo}}(POLYNÔMES COMPLEXES)

\fbox{\textbf{\exo}} (POLYNÔMES COMPLEXES)

\fbox{\textbf{\exo}} (POLYNÔMES COMPLEXES)

\fbox{\textbf{\exo}}

\end{document}
